% 3.threat_model.tex — Threat Model

\section{Threat Model}
\label{sec:threat_model}

As emotion-aware ALLMs enter deployment in affective companion systems, mental health support, and intelligent customer service, their downstream responses become conditioned on an internal emotion percept---the model's implicit interpretation of the speaker's emotional state---that the user can neither observe nor correct.
Section~\ref{sec:obs_q2} demonstrated that even a single misperceived emotion is enough to trigger systematic response degradation---faithfulness collapse on neutral inputs, hallucinated emotional narratives, and a pronounced vulnerability gap between positive and negative valences.
This fragility motivates a concrete adversarial question: can an attacker reliably steer the model's emotion percept to a chosen target, using only an imperceptible audio perturbation?

\paragraph{Attacker's goal.}
Given clean audio $x$ with ground-truth speaker emotion $e_s$, the attacker seeks a perturbation $\delta$ such that the adversarial audio $x' = x + \delta$ causes the target ALLM to output a pre-specified emotion $e_t \neq e_s$, while the linguistic content and perceptual quality of the waveform remain intact.
The attack is \textit{targeted}: the attacker controls not only \textit{whether} the prediction changes but \textit{which} emotion it changes to---a strictly harder objective than untargeted disruption.

\paragraph{Attacker's knowledge.}
We operate in a white-box setting: the attacker has full access to the target ALLM's architecture, weights, and intermediate representations, and can compute end-to-end gradients through the entire audio-encoder--adapter--LLM pipeline.
The attacker knows the prompt templates used for emotion elicitation but cannot modify model parameters or prompts; only the input waveform is under adversarial control.
We assume the adversarial audio is injected as a digital input (e.g., uploaded file or API-mediated stream) rather than played over-the-air, so channel distortions are not modeled.

\paragraph{Constraints.}
The perturbation must satisfy three requirements.
An $L_\infty$ norm bound $\|\delta\|_\infty \leq \varepsilon$ with $\varepsilon = 0.008$ enforces bounded sample-level waveform perturbation as a hard constraint, applied by projection after each optimization step.
Semantic preservation requires that the model's own transcription of the adversarial audio remains highly similar to that of the clean audio, ensuring the linguistic content is unaltered.
A multi-resolution STFT magnitude loss further penalizes perceptually salient spectral distortions.
The latter two are implemented as soft loss terms whose formulations are given in Section~\ref{sec:methodology}.

\medskip\noindent
\textbf{Formal problem statement.}
Let $f$ denote the target ALLM, $p_{\text{emo}}$ an emotion elicitation prompt, and $p_{\text{asr}}$ a transcription prompt.
Given clean audio $x \in \mathbb{R}^T$ with source emotion $e_s$ and an attacker-chosen target $e_t \neq e_s$, the objective is to find $\delta$ such that the adversarial audio $x' = x + \delta$ satisfies:
%
\begin{enumerate}
    \item[\textnormal{(i)}] $f(x',\, p_{\text{emo}}) = e_t$ \hfill \textit{(targeted emotion flip)}
    \item[\textnormal{(ii)}] $\mathrm{sim}\!\bigl(\mathrm{ASR}_{f}(x',\, p_{\text{asr}}),\; \mathrm{ASR}_{f}(x,\, p_{\text{asr}})\bigr) \geq \tau$ \hfill \textit{(semantic preservation)}
    \item[\textnormal{(iii)}] $\|\delta\|_\infty \leq \varepsilon$ \hfill \textit{(imperceptibility)}
\end{enumerate}
%
where $\mathrm{ASR}_{f}(\cdot,\, p_{\text{asr}})$ denotes the model's transcription output and $\mathrm{sim}(\cdot,\,\cdot)$ is a semantic similarity measure detailed in Section~\ref{sec:methodology}.
